\begin{examplebox}
\textbf{例 1 (直接套用, 同时检验条件)}

椭圆 $\frac{x^2}{4}+y^2=1$, 在切点 $P$ 有 $r_1r_2=2$.
求中心到 $P$ 处切线的距离 $d$.

\textit{解}: $a=2,b=1$, 因此
\[
d=\frac{ab}{\sqrt{r_1r_2}}=\frac2{\sqrt2}=\boxed{\sqrt2}.
\]
\textit{检验}: $b^2=1\le2\le4=a^2$, 已知乘积确实可能出现.

\medskip
\textbf{例 2 (最值, 修正原题错误)}

求椭圆 $\frac{x^2}{a^2}+\frac{y^2}{b^2}=1$ 上的切线到中心距离 $d$ 的最大值、最小值及对应切点.

\textit{解法 A, 一眼判断}: 在两轴端点处, 切线分别为 $x=\pm a$ 和 $y=\pm b$, 对应距离是 $a$ 和 $b$.
但必须证明其他点的距离不会超出此范围.

\textit{解法 B, 严格最值}: 由焦半径公式
\[
r_1r_2=a^2-\frac{c^2x_0^2}{a^2}.
\]
因为 $0\le x_0^2\le a^2$,
\[
b^2=a^2-c^2\le r_1r_2\le a^2.
\]
把上述不等式代入 $d=ab/\sqrt{r_1r_2}$, 注意取倒数会反向:
\[
\boxed{b\le d\le a}.
\]
$d_{\max}=a$, 当且仅当 $x_0=\pm a,y_0=0$, 即\textbf{长轴顶点};
$d_{\min}=b$, 当且仅当 $x_0=0,y_0=\pm b$, 即\textbf{短轴顶点}.
\textit{易错提醒}: 不能把 ``短轴顶点处最大'' 写成题目结论.

\medskip
\textbf{例 3 (只给焦半径差, 不必设切点)}

椭圆 $\frac{x^2}{25}+\frac{y^2}{9}=1$ 上某切点满足 $|r_1-r_2|=6$.
求中心到该切线的距离.

\textit{解}: 由定义 $r_1+r_2=2a=10$, 两数的乘积可用平方差求出:
\[
4r_1r_2=(r_1+r_2)^2-(r_1-r_2)^2
=100-36=64.
\]
于是 $r_1r_2=16$,
\[
\boxed{d=\frac{5\times3}{\sqrt{16}}=\frac{15}{4}}.
\]
\textit{可行性核查}: $|r_1-r_2|=6\le2c=8$, 数据合理.

\medskip
\textbf{例 4 (反求所有切点, 四个对称解)}

椭圆 $\frac{x^2}{25}+\frac{y^2}{9}=1$ 的某条切线满足 $d=15/4$.
求切点 $P$.

\textit{解}: $ab=15$, 故
\[
r_1r_2=\left(\frac{ab}{d}\right)^2=16.
\]
又 $a=5,c=4,e=4/5$,
\[
16=25-\frac{16}{25}x_0^2
\quad\Longrightarrow\quad x_0^2=\frac{225}{16}.
\]
所以 $x_0=\pm15/4$. 再代入椭圆方程,
\[
\frac{y_0^2}{9}=1-\frac{(225/16)}{25}
=\frac7{16},\qquad y_0=\pm\frac{3\sqrt7}{4}.
\]
\textbf{四个切点均满足}:
\[
\boxed{P=\left(\pm\frac{15}{4},\ \pm\frac{3\sqrt7}{4}\right)}
\]
(两个正负号独立选择, 共四个点).

\medskip
\textbf{例 5 (两种方法互相验算)}

椭圆 $\frac{x^2}{25}+\frac{y^2}{9}=1$ 上取
$P(3,12/5)$, 求中心到此处切线的距离.

\textit{解法 A, 利用 C012}:
\[
r_1r_2=25-\frac{16}{25}\cdot 9=\frac{481}{25},
\qquad
\boxed{d=\frac{15}{\sqrt{481}/5}=\frac{75}{\sqrt{481}}}.
\]
\textit{解法 B, 直接写切线}:
\[
\frac{3x}{25}+\frac{4y}{15}=1,
\qquad
 d=\frac1{\sqrt{9/625+16/225}}
=\frac{75}{\sqrt{481}}.
\]
两种方法相互吻合. 其中 $P$ 确实在椭圆上:
$9/25+(144/25)/9=9/25+16/25=1$.

\medskip
\textbf{例 6 (几何角度: 切线距离与焦点三角形顶角)}

已知椭圆半长轴 $a=5$, 半短轴 $b=3$,
且切点处的焦点三角形顶角满足 $\cos(\varphi/2)=4/5$.
求 $d$ 以及 $r_1r_2$.

\textit{解}: 几何公式给出
\[
d=a\cos\frac\varphi2=5\cdot\frac45=\boxed{4}.
\]
再由 $r_1r_2\cos^2(\varphi/2)=b^2$,
\[
r_1r_2=\frac9{(4/5)^2}=\boxed{\frac{225}{16}}.
\]
注意此例的 $d=4$ 在 $[b,a]=[3,5]$ 内, 且焦半径乘积
$225/16$ 在 $[9,25]$ 内, 两项均符合范围检验.
\end{examplebox}
